\chapter{Arthroscopy of the Shoulder}
\index{Arthroscopy of the Shoulder}

\section*{Definition and Overview}
Arthroscopy of the shoulder is a keyhole operation in which a small camera and surgical instruments are passed through tiny incisions to diagnose and treat problems within and around the shoulder joint. It is commonly used to repair torn rotator cuff tendons, relieve subacromial impingement, treat instability and dislocation, repair labral tears, or release a frozen shoulder. Because the shoulder is a shallow ball-and-socket joint surrounded by many tendons, ligaments, and bursae, arthroscopy allows the surgeon to inspect and treat the whole joint without a large incision. Most procedures are performed as day surgery under general anaesthesia, often with a nerve block.

\section*{Epidemiology}
Shoulder problems are very common, and arthroscopy of the shoulder has become one of the most frequent orthopaedic procedures. Rotator cuff tears become increasingly common with age, affecting a substantial proportion of older adults, although many are painless and never require surgery. Shoulder dislocations and recurrent instability are most common in young, active people, particularly men. Subacromial impingement and frozen shoulder are also frequent reasons for shoulder surgery, and the rising participation in overhead and throwing sports has increased the demand for arthroscopic treatment.

\section*{Aetiology and Pathophysiology}
\begin{itemize}
    \item \textbf{Rotator cuff tears:} tears of the tendons around the shoulder, arising from acute injury or from age-related degeneration of the tendon
    \item \textbf{Impingement:} compression of the tendons and bursa beneath the bony acromion, causing pain with overhead movement
    \item \textbf{Instability and dislocation:} damage to the labrum and capsuloligamentous structures when the shoulder slips or dislocates, allowing the humeral head to displace
    \item \textbf{Labral tears:} tears of the rim of cartilage that deepens the socket, sometimes associated with injury to the long head of the biceps tendon
    \item \textbf{Frozen shoulder (adhesive capsulitis):} thickening and contraction of the joint capsule causing global stiffness
    \item \textbf{Loose bodies, synovitis, and joint infection:} less common indications for arthroscopic treatment
\end{itemize}

\section*{Clinical Features}
\begin{itemize}
    \item Shoulder pain, particularly at night or with overhead activities
    \item Weakness or difficulty lifting the arm against resistance
    \item A painful arc of movement during the middle of lifting the arm
    \item Clicking, catching, or grinding sensations within the shoulder
    \item A feeling that the shoulder is loose, or repeated episodes of dislocation
    \item Stiffness and restricted range of movement
    \item After surgery: mild pain, swelling, and bruising around the incisions for a few days
\end{itemize}

\section*{Differential Diagnosis}
Conditions that can mimic the symptoms treated by shoulder arthroscopy include:
\begin{itemize}
    \item Cervical spine problems, such as a nerve root compression in the neck referring pain to the shoulder
    \item Osteoarthritis of the glenohumeral or acromioclavicular joints
    \item Frozen shoulder without an underlying tendon problem
    \item Calcific tendinopathy of the rotator cuff
    \item Subacromial bursitis and other causes of subacromial pain
    \item Rarely, referred pain from the heart, diaphragm, or lung
\end{itemize}

\section*{When to Seek Medical Care}
See a doctor if shoulder pain or weakness persists for more than a few weeks, is worse at night, or interferes with daily activities and sleep.

\textbf{Contact your local emergency services (or attend an emergency department) for:}
\begin{itemize}
    \item A shoulder that is obviously deformed, dislocated, or has not gone back into place
    \item Sudden severe shoulder pain after an injury or fall
    \item Shoulder pain accompanied by chest pain, breathlessness, or collapse
    \item After arthroscopy: fever, spreading redness, or severe pain not controlled by prescribed analgesia
\end{itemize}

\section*{Diagnosis and Investigations}
\begin{itemize}
    \item \textbf{History:} how the symptoms started, the activities that aggravate them, and any previous dislocations
    \item \textbf{Examination:} assessment of range of movement, strength testing, impingement signs, and instability provocation tests
    \item \textbf{X-ray:} to exclude arthritis, calcification, or signs of previous dislocation
    \item \textbf{Ultrasound:} dynamic assessment of the rotator cuff tendons and any fluid
    \item \textbf{MRI or MR arthrogram:} the most accurate imaging for tendon, labral, and ligament tears
    \item \textbf{Arthroscopy:} direct visual inspection of the joint, confirming the diagnosis and allowing treatment in the same procedure
\end{itemize}

\section*{Management}
\begin{itemize}
    \item \textbf{Subacromial decompression:} removal of inflamed bursa and trimming of bone spurs to relieve impingement
    \item \textbf{Rotator cuff repair:} stitching of the torn tendon back to its bony attachment, usually with suture anchors
    \item \textbf{Labral repair or stabilisation:} reattachment of the torn labrum and tightening of the capsule and ligaments to prevent further dislocations
    \item \textbf{Capsular release:} division of the contracted capsule to treat frozen shoulder
    \item \textbf{Loose body removal and washout:} for loose fragments, synovitis, or infection
    \item \textbf{Sling use:} for a period of time after tendon or labral repair to protect the healing tissue
    \item \textbf{Physiotherapy:} a graded rehabilitation programme, which is essential after most shoulder arthroscopy procedures
\end{itemize}

\section*{Complications}
\begin{itemize}
    \item Stiffness, which is the most common problem after rotator cuff or labral repair
    \item Failure of the repair to heal, with recurrent weakness, pain, or instability
    \item Infection in the joint or at the incisions, which is uncommon
    \item Nerve damage, most often temporary numbness or weakness
    \item Ongoing pain, or an inability to return fully to sport
    \item Anaesthetic-related risks, as with any operation
\end{itemize}

\section*{Prognosis and Outcomes}
Results are generally good. Most patients obtain substantial relief from pain and improvement in function, although recovery is often slow: rotator cuff repairs commonly require months of rehabilitation, with full strength returning at around nine to twelve months. Stabilisation surgery prevents further dislocations in the majority of cases. Some residual stiffness or discomfort may persist, and the final outcome depends on the extent of the damage, the quality of the repair, and commitment to rehabilitation.

\section*{Prevention and Self-Care}
\begin{itemize}
    \item Strengthen the shoulder, neck, and upper back muscles regularly
    \item Use correct technique and allow adequate recovery in overhead and throwing sports
    \item Warm up before activity and increase load gradually
    \item Maintain good posture, especially during desk work
    \item Follow the rehabilitation exercises exactly after surgery
    \item Seek early review of persistent shoulder pain rather than waiting for it to become severe
\end{itemize}

\section*{Sources and Further Reading}
The following organisations provide reliable information on shoulder arthroscopy for patients and healthcare students.

\begin{itemize}
    \item NHS. \textit{Arthroscopy of the shoulder: overview and recovery.} https://www.nhs.uk/conditions/arthroscopy/
    \item American Academy of Orthopaedic Surgeons (AAOS). \textit{Shoulder arthroscopy.} https://orthoinfo.aaos.org/en/treatment/shoulder-arthroscopy/
    \item Mayo Clinic. \textit{Arthroscopy overview.} https://www.mayoclinic.org/tests-procedures/arthroscopy/
    \item MedlinePlus. \textit{Shoulder arthroscopy.} https://medlineplus.gov/ency/article/007199.htm
    \item Versus Arthritis. \textit{Shoulder pain information for patients.} https://www.versusarthritis.org/about-arthritis/conditions/shoulder-pain/
    \item MSD Manual. \textit{Shoulder pain and rotator cuff disorders.} https://www.msdmanuals.com/professional/musculoskeletal-and-connective-tissue-disorders/approach-to-the-patient-with-joint-disease/shoulder-pain
\end{itemize}
